\section{Automated Code-Review Comment Generation}
\label{sec:sota-generation}

The dominant framing treats review-comment generation as transduction from a code
change to a natural-language comment. CodeReviewer pre-trains a model on a large
multilingual corpus of review data and fine-tunes it for comment generation,
quality estimation and code refinement~\cite{li2022codereviewer}; Tufano et al.\
establish that pre-trained models outperform earlier task-specific
sequence-to-sequence systems on the same
task~\cite{tufano2022codereview}. Instruction-following general-purpose models
subsequently made fine-tuning optional for many of these tasks, a shift documented
in two systematic reviews of language models in software
engineering~\cite{hou2023llmse,fan2023llmse}.

The unresolved question in this line is what the model is given. In each of the
systems above, the input is the change, usually a diff and sometimes the
enclosing function. The model's knowledge of everything else comes from
pre-training. Recent benchmark work has begun to treat this as the limiting
factor rather than a detail of formatting. Hu et al.'s ContextCRBench is built
specifically to test fine-grained review with enriched context and argues that
textual context alone constrains what models can
achieve~\cite{hu2025contextcrbench}. CodeFuse-CR-Bench evaluates review at
repository level across seventy projects, on the premise that comprehensiveness
cannot be assessed from a diff in
isolation~\cite{guo2025codefusecrbench}. SWR-Bench assembles a thousand manually
verified pull requests for real-world review-comment
generation~\cite{zeng2025swrbench}.

These benchmarks establish that context matters and supply data on which to
measure it. What they do not do is hold the model, the prompt and the evaluation
fixed while varying \emph{how} the context is selected. Enriching context and
comparing selection strategies are different experiments, and it is the second that
this thesis runs.

Adjacent work targets other pull-request artefacts. Mariotto et al.\ score pull
requests against engineering competencies at scale and show that prompt design
substantially improves generated pull-request \emph{descriptions} relative to naive
baselines~\cite{mariotto2025esem,mariotto2025}; the dataset assembled for that work
is the origin of the pull requests used here, as
Section~\ref{sec:dataset} records. The distinction is that descriptions summarise a
change the author already understands, whereas review comments must reason about
consequences the author may have missed. Out-of-diff context should therefore
matter more for the latter.
